% Part: first-order-logic
% Chapter: tableaux
% Section: derivations

\documentclass[../../../include/open-logic-section]{subfiles}

\begin{document}

\iftag{FOL}
      {\olfileid{fol}{tab}{der}}
      {\olfileid{pl}{tab}{der}}

\olsection{\usetoken{P}{tableau}}

\begin{explain}
ધારણા શું છે તે આપણે કહી દીધું છે અને અનુમાનના નિયમો આપ્યા છે. આમાંથી
!!^{tableau}sનું આગમનાત્મક રીતે નિર્માણ થાય છે: દરેક !!{tableau} કાં તો
એક કે વધુ ધારણાઓથી બનેલી એક જ શાખા છે, અથવા કોઈ !!a{tableau}ની શાખા
પર અનુમાનનો કોઈ એક નિયમ લાગુ કરવાથી મળે છે.
\end{explain}

\begin{defn}[!!^{tableau}]
ધારણાઓ $\sFmla{S_1}{!A_1}$, \dots, $\sFmla{S_n}{!A_n}$ માટેનો
!!^a{tableau} (જ્યાં દરેક $S_i$ કાં તો $\True$ અથવા~$\False$ છે) એ
નીચેની શરતો સંતોષતું !!{signed formula}sનું સાન્ત વૃક્ષ છે:
\begin{enumerate}
\item વૃક્ષનાં સૌથી ઉપરનાં $n$ !!{signed formula}s એકબીજાની નીચે
  $\sFmla{S_i}{!A_i}$ છે.
\item વૃક્ષમાં જે દરેક !!{signed formula} ધારણાઓમાંથી એક નથી, તે તેની
  ઉપરની શાખામાંના !!a{signed formula} પર અનુમાન-નિયમના યોગ્ય ઉપયોગથી
  મળે છે.
\end{enumerate}
!!a{tableau}ની કોઈ શાખા \emph{બંધ} છે જો અને માત્ર જો તેમાં
$\sFmla{\True}{!A}$ અને~$\sFmla{\False}{!A}$ બંને હોય; અન્યથા તે
\emph{ખુલ્લી} છે. !!^a{tableau}ની દરેક શાખા બંધ હોય તો તે (તેની
ધારણાઓના ગણ માટેનો) \emph{બંધ !!{tableau}} છે. જો !!a{tableau} બંધ ન
હોય, એટલે કે તેમાં ઓછામાં ઓછી એક ખુલ્લી શાખા હોય, તો તે \emph{ખુલ્લો}
છે.
\end{defn}

\begin{ex}
  ધારણાઓનો દરેક ગણ પોતે જ !!a{tableau} છે, પરંતુ સામાન્ય રીતે તે બંધ
  નહીં હોય. (દેખીતી રીતે, ધારણાઓમાં !!{signed formula}sની
  $\sFmla{\True}{!A}$ અને~$\sFmla{\False}{!A}$ જોડી પહેલેથી હોય તો જ
  તે બંધ છે.)

  કોઈ !!a{tableau}માંના !!a{signed formula}~$!A$ને અનુમાનનો કોઈ એક
  નિયમ લાગુ કરીને તેમાંથી નવો અને મોટો ટેબ્લો મેળવી શકીએ છીએ, પછી ભલે
  પહેલો ટેબ્લો ખુલ્લો હોય કે બંધ. જે~$!A$ના આવર્તનને નિયમ લાગુ કરીએ તે
  આવર્તન ધરાવતી દરેક શાખાના છેડે નિયમ એક કે વધુ !!{signed formula}s
  ઉમેરશે.

  દાખલા તરીકે, ધારણા $\sFmla{\True}{!A \land \lnot !A}$ લો. માત્ર એ
  ધારણાથી બનેલો (ખુલ્લો) !!{tableau} આ રહ્યો:
  \begin{center}
    \begin{tableau}{}
      [\sFmla{\True}{\formula{A} \land \lnot \formula{A}}, just=\TAss]
    \end{tableau}{}
  \end{center}
  ધારણાને $\TRule{\True}{\land}$ નિયમ લાગુ કરીને તેમાંથી નવો
  !!{tableau} મળે છે. એ નિયમ !!{tableau}માં બે નવી પંક્તિઓ,
  $\sFmla{\True}{!A}$ અને $\sFmla{\True}{\lnot !A}$, ઉમેરવાની છૂટ આપે છે:
  \begin{center}
    \begin{tableau}{}
      [\sFmla{\True}{\formula{A} \land \lnot \formula{A}}, just=\TAss
        [\sFmla{\True}{\formula{A}}, just={\TRule{\True}{\land}[1]},
          [\sFmla{\True}{\lnot \formula{A}}, just={\TRule{\True}{\land}[1]}
          ]
        ]
      ]
    \end{tableau}{}
  \end{center}
  !!{tableau}s લખતી વખતે આપણે લાગુ કરેલા નિયમો જમણી બાજુ નોંધીએ છીએ
  (દાખલા તરીકે, $\TRule{\True}{\land} 1$નો અર્થ છે કે તે પંક્તિ પરનું
  !!{signed formula}, પંક્તિ~$1$ પરના !!{signed formula}ને
  $\TRule{\True}{\land}$ નિયમ લાગુ કરવાથી મળ્યું છે). હવે આ નવા
  !!{tableau}માં વધારાનાં !!{signed formula}s છે, પરંતુ તેમાંથી માત્ર
  એકને ($\sFmla{\True}{\lnot !A}$ને) નિયમ લાગુ કરી શકાય છે (આ કિસ્સામાં
  $\TRule{\True}{\lnot}$ નિયમ). તેનાથી નીચેનો બંધ !!{tableau} મળે છે:
  \begin{center}
    \begin{tableau}{}
      [\sFmla{\True}{\formula{A} \land \lnot \formula{A}}, just=\TAss
        [\sFmla{\True}{\formula{A}}, just={\TRule{\True}{\land}[1]}
          [\sFmla{\True}{\lnot \formula{A}}, just={\TRule{\True}{\land}[1]}
            [\sFmla{\False}{\formula{A}}, just={\TRule{\True}{\lnot}[3]}, close]
          ]
        ]
      ]
    \end{tableau}{}
  \end{center}
\end{ex}

\end{document}
